\section{체확장과 차수}
\label{sec:field-degree}

\begin{learninggoal}
체확장을 벡터공간으로 해석하고 차수를 계산한다. 중간체를 포함하는 확장탑에서 탑 법칙을 증명하고 활용한다.
\end{learninggoal}

\subsection{체확장은 벡터공간이다}

\begin{definition}[체확장]
두 체 $K,L$에 대해 $K$가 $L$의 부분체이면 $L/K$를 \emph{체확장}(field extension)이라 한다. $K\subseteq E\subseteq L$인 부분체 $E$를 $L/K$의 \emph{중간체}(intermediate field)라 한다.
\end{definition}

$K\subseteq L$이면 $L$의 덧셈과 $K$의 원소에 의한 곱셈을 이용해 $L$을 $K$-벡터공간으로 볼 수 있다.

\begin{definition}[확장의 차수]
체확장 $L/K$의 차수를
\[
  [L:K]=\dim_K L
\]
로 정의한다. 이 차원이 유한하면 $L/K$를 \emph{유한확장}, 그렇지 않으면 \emph{무한확장}이라 한다.
\end{definition}

\begin{example}
\begin{enumerate}[label=(\alph*)]
  \item $\C/\R$에서 $1,i$는 기저이므로 $[\C:\R]=2$이다.
  \item $\Q(\sqrt2)/\Q$에서 $1,\sqrt2$는 기저이므로 차수는 $2$이다.
  \item $\R/\Q$는 무한확장이다.
  \item 유리함수체 $K(t)/K$는 무한확장이다. 실제로 $1,t,t^2,\dots$는 $K$ 위에서 선형독립이다.
\end{enumerate}
\end{example}

\begin{pitfall}
$[L:K]$는 집합의 원소 수가 아니다. 무한체 $K$에 대해서도 유한차수 확장이 가능하며, $[\C:\R]=2$이지만 두 체 모두 무한집합이다.
\end{pitfall}

\subsection{탑 법칙}

\begin{theorem}[탑 법칙]
체확장탑
\[
  K\subseteq L\subseteq M
\]
에 대해
\[
  [M:K]=[M:L][L:K]
\]
가 성립한다. 유한차수인 경우 오른쪽은 보통의 자연수 곱이고, 어느 한 차수가 무한이면 왼쪽도 무한이다.
\end{theorem}

\begin{proofidea}
$L/K$의 기저와 $M/L$의 기저를 하나씩 곱하면 $M/K$의 기저가 된다.
\end{proofidea}

\begin{proof}
먼저 $[L:K]=m$, $[M:L]=n$이 유한하다고 하자. $L/K$의 기저를
\[
  x_1,\dots,x_m
\]
이라 하고 $M/L$의 기저를
\[
  y_1,\dots,y_n
\]
이라 하자. 그러면
\[
  \{x_i y_j:1\le i\le m,\ 1\le j\le n\}
\]
가 $M/K$의 기저임을 보이면 된다.

먼저 생성성을 보자. 임의의 $z\in M$는
\[
  z=\sum_{j=1}^n c_jy_j
  \qquad(c_j\in L)
\]
로 쓰이고, 각 $c_j$는
\[
  c_j=\sum_{i=1}^m a_{ij}x_i
  \qquad(a_{ij}\in K)
\]
로 쓰인다. 따라서
\[
  z=\sum_{i,j}a_{ij}x_i y_j
\]
이다.

선형독립성을 보자. $a_{ij}\in K$에 대해
\[
  \sum_{i,j}a_{ij}x_i y_j=0
\]
이라 하자. $y_j$에 대한 $L$-선형결합으로 묶으면
\[
  \sum_j\left(\sum_i a_{ij}x_i\right)y_j=0.
\]
$y_j$들의 $L$-선형독립성에서 각 $j$에 대해
\[
  \sum_i a_{ij}x_i=0
\]
이고, 다시 $x_i$들의 $K$-선형독립성에서 모든 $a_{ij}=0$이다. 따라서 곱으로 만든 $mn$개의 원소가 기저이다.

이제 $[L:K]$ 또는 $[M:L]$이 무한이라고 하자. 임의로 큰 유한 선형독립 집합을 각각 택하고 위와 같은 곱셈 논증을 적용하면 $M$ 안에 임의로 큰 $K$-선형독립 집합이 생긴다. 따라서 $[M:K]$는 무한이다.
\end{proof}

\begin{corollary}
$L/K$가 유한확장이고 $E$가 중간체이면
\[
  [E:K]\mid[L:K],
  \qquad
  [L:E]\mid[L:K].
\]
\end{corollary}

\begin{corollary}[소수 차수 확장]
$[L:K]=p$가 소수이면 중간체는 $K$와 $L$뿐이다.
\end{corollary}

\begin{proof}
중간체 $E$에 대해
\[
  p=[L:E][E:K]
\]
이다. 두 인수는 양의 정수이므로 하나가 $1$이어야 한다.
\end{proof}

\begin{example}
$\Q\subseteq\Q(\sqrt2)\subseteq\Q(\sqrt2,\sqrt3)$를 생각하자. 뒤에서
\[
  [\Q(\sqrt2,\sqrt3):\Q(\sqrt2)]=2
\]
임을 확인할 것이다. 그러면 탑 법칙으로
\[
  [\Q(\sqrt2,\sqrt3):\Q]=4
\]
이다.
\end{example}

\subsection{유한체의 원소 수}

\begin{proposition}
유한체 $F$의 표수가 $p$이면 어떤 양의 정수 $n$에 대해
\[
  |F|=p^n
\]
이다.
\end{proposition}

\begin{proof}
$F$의 소체는 $\F_p$이고, $F$는 유한집합이므로 $\F_p$ 위에서 유한차원이다. $n=[F:\F_p]$라 하면 $F$는 $n$차원 벡터공간이다. 각 기저좌표에는 $p$개의 선택이 있으므로 원소 수는 $p^n$이다.
\end{proof}

\begin{corollary}
원소가 $6$개인 체는 존재하지 않는다.
\end{corollary}

\begin{checkpoint}
\begin{enumerate}[label=(\alph*)]
  \item $[\Q(\sqrt2):\Q]=2$임을 기저로 증명하라.
  \item $[M:K]=15$인 유한확장의 중간체 $E$가 가질 수 있는 차수를 모두 적어라.
  \item 원소가 $12$개인 체가 존재할 수 있는지 판정하라.
\end{enumerate}
\end{checkpoint}
